\documentclass[10pt,a4paper]{amsart}
\usepackage[margin=19mm]{geometry}
\usepackage{amsmath,amssymb,mathtools}
\usepackage[scr=rsfs,cal=euler]{mathalfa}
\usepackage{xfrac}
\usepackage[unicode,hidelinks,bookmarks=false]{hyperref}
\newcommand{\ie}{i.e.\ }
\newcommand{\cf}{cf.\ }
\newcommand{\resp}{respectively\ }
\newcommand{\Subsection}{\subsection}
\providecommand{\nobreakdash}{\nobreak\hbox{-}}
\setlength{\emergencystretch}{2em}
\multlinegap0pt
\usepackage{amssymb}
\usepackage{xcolor}
\usepackage{algorithm}\usepackage{algpseudocode}
\newtheorem{proposition}{Proposition}
\title{Hodge Spectral Surrogates for Topology-Constrained Optimization}
\author{Satoshi Kanno, Yoshi-aki Shimada}
\date{Source: arXiv:2606.25194v2 (CC BY 4.0)}
\begin{document}
\maketitle

\noindent\emph{Source and licence.} Hodge Spectral Surrogates for Topology-Constrained Optimization. Version arXiv:2606.25194v2 (CC BY 4.0), distributed by arXiv under the Creative Commons Attribution 4.0 International licence, http://creativecommons.org/licenses/by/4.0/, as recorded in the arXiv licence metadata for this exact version and shown on the abstract page https://arxiv.org/abs/2606.25194v2. Full licence notice: \texttt{licenses/arxiv-2606.25194-CC-BY-4.0.txt}.\par\smallskip

\noindent\textit{Editorial scope.} Counted unit: the article's proposition proposition (source0.tex lines 1917--1989). The article's complete original proof of it is delivered with it (source0.tex lines 1991--2146, one \texttt{proof} environment of 156 lines). No mathematics is written, repaired or re-derived: the statement and the proof are the article's own lines, in the article's own notation, and no retained line is substituted or re-flowed.  No further source-local context is needed: every cross-reference of every retained line resolves inside the two delivered spans.  Nothing was omitted from any retained span, so the package carries no omission record.  No retained line contains a citation, so the wrapper carries no bibliography and no bibliography entry.  No retained line contains a figure environment, an image-inclusion command or drawing code, so no image is delivered and the package declares no asset dependency.  Editorial formatting only, in addition: an AMS \texttt{amsart} preamble that restates the article's own declarations and macros used by the retained lines, namely the environments \texttt{proposition} on separate counters, together with the standard packages those lines need.  The retained environments print in this document as Proposition 1 in order of appearance, whereas the article prints the same items with the numbering of its own full document.  The complete native source of the article as distributed in the arXiv e-print for this exact version is bundled unchanged as \texttt{sources/203585/source0.tex}.
\begin{proposition}[Hard-limit consistency of the \(W_q\)-weighted penalty-free moment]
Let \(K\) be the corresponding hard complex and define 
$n_q
=
|K^{(q)}|,
~
b
=
\beta_q(K)$.
With respect to the decomposition
\[
C_q^{\mathrm{amb}}
=
C_q(K)
\oplus
C_q(K)^\perp,
\]
the hard limit of the penalty-free operator satisfies
\[
L_{q,\mathrm{wc}}^{\mathrm{hard}}
=
L_q(K)
\oplus
0,
\]
whereas
\[
W_q^{\mathrm{hard}}
=
I_{n_q}
\oplus
0.
\]
Consequently,
\[
\overline S_{q,d}^{W,\mathrm{hard}}
=
\frac{1}{n_q}
\operatorname{Tr}_{C_q(K)}
\left(
I-
\frac{
L_q(K)
}{
\Lambda_q^{\mathrm{amb}}
}
\right)^d.
\]
If \(b<n_q\), let \(\gamma_q>0\) be the smallest positive eigenvalue of \(L_q(K)\). Then
\[
0
\le
\overline S_{q,d}^{W,\mathrm{hard}}
-
\overline\beta_q(K)
\le
\left(
1-
\frac{
\gamma_q
}{
\Lambda_q^{\mathrm{amb}}
}
\right)^d.
\]
If \(b=n_q\), then
$\overline S_{q,d}^{W,\mathrm{hard}}
=
\overline\beta_q(K)
=
1$
for every \(d\).
\end{proposition}

\begin{proof}
In the hard limit,
\[
R_p^{(\varepsilon)}
\longrightarrow
\Pi_p.
\]
The projected boundary operators coincide with the ordinary boundary operators of \(K\) on the active chain spaces and vanish on inactive simplex directions. Hence
\[
L_{q,\mathrm{wc}}^{\mathrm{hard}}
=
L_q(K)\oplus 0.
\]
At the same time,
\[
W_q^{\mathrm{hard}}
=
\Pi_q
=
I_{n_q}\oplus 0.
\]
Therefore,
\[
W_q^{\mathrm{hard}}
\left(
I-
\frac{
L_{q,\mathrm{wc}}^{\mathrm{hard}}
}{
\Lambda_q^{\mathrm{amb}}
}
\right)^d
=
\left(
I-
\frac{
L_q(K)
}{
\Lambda_q^{\mathrm{amb}}
}
\right)^d
\oplus
0,
\]
which proves
\[
\overline S_{q,d}^{W,\mathrm{hard}}
=
\frac{1}{n_q}
\operatorname{Tr}_{C_q(K)}
\left(
I-
\frac{
L_q(K)
}{
\Lambda_q^{\mathrm{amb}}
}
\right)^d.
\]

Now suppose \(b<n_q\), and let the eigenvalues of \(L_q(K)\) be ordered as
\[
0
=
\lambda_1
=
\cdots
=
\lambda_b
<
\lambda_{b+1}
\le
\cdots
\le
\lambda_{n_q}.
\]
Since the zero eigenvalues contribute one to the polynomial filter,
\[
\overline S_{q,d}^{W,\mathrm{hard}}
=
\frac{b}{n_q}
+
\frac{1}{n_q}
\sum_{i>b}
\left(
1-
\frac{
\lambda_i
}{
\Lambda_q^{\mathrm{amb}}
}
\right)^d.
\]
For every positive eigenvalue,
\[
\lambda_i\ge\gamma_q.
\]
Hence,
\[
0
\le
\left(
1-
\frac{
\lambda_i
}{
\Lambda_q^{\mathrm{amb}}
}
\right)^d
\le
\left(
1-
\frac{
\gamma_q
}{
\Lambda_q^{\mathrm{amb}}
}
\right)^d.
\]
It follows that
\[
\begin{aligned}
0
&\le
\overline S_{q,d}^{W,\mathrm{hard}}
-
\frac{b}{n_q}
\\
&\le
\frac{
n_q-b
}{
n_q
}
\left(
1-
\frac{
\gamma_q
}{
\Lambda_q^{\mathrm{amb}}
}
\right)^d
\\
&\le
\left(
1-
\frac{
\gamma_q
}{
\Lambda_q^{\mathrm{amb}}
}
\right)^d.
\end{aligned}
\]
This proves the result.
\end{proof}

\end{document}
